% ---- begin paper/tables/t2_largo_plazo.tex
\begin{table}[htbp]\centering
\begin{minipage}{460pt}
\caption{The impact of the earthquake on cognitive and non-cognitive development in the medium- and long-term.}
\label{tab:longrun}
{\small
\begin{tabular*}{460pt}{@{}l@{\extracolsep{\fill}}llll@{}}
\toprule
 & \multicolumn{2}{@{}l}{Medium-term effect (Gillmore)} & \multicolumn{2}{@{}l}{Long-term effect} \\
 & \multicolumn{2}{@{}l}{2017 wave} & \multicolumn{2}{@{}l}{2024 wave} \\
\cmidrule(r){2-3}\cmidrule(l){4-5}
 & (1) & (2) & (3) & (4) \\
\midrule
\textbf{\textit{Panel A: Peabody}} & & & & \\
Affected$\times$Earthquake & $-$0.174\sym{**} & $-$0.280\sym{***} & $-$0.240\sym{**} & $-$0.412\sym{**} \\
 & (0.0686) & (0.0866) & (0.104) & (0.180) \\
\addlinespace[9pt]
Observations & 14,069 & 14,069 & 13,779 & 13,779 \\
\addlinespace[9pt]
\textbf{\textit{Panel B: CBCL2}} & & & & \\
Affected$\times$Earthquake & $-$0.129\sym{*} & $-$0.220\sym{**} & 0.084 & $-$0.066 \\
 & (0.0763) & (0.105) & (0.0884) & (0.125) \\
\addlinespace[9pt]
Observations & 11,568 & 11,568 & 11,038 & 11,038 \\
\addlinespace[9pt]
Municipality FE & \checkmark & \checkmark & \checkmark & \checkmark \\
Birth Cohort FE & \checkmark & \checkmark & \checkmark & \checkmark \\
Child-Mother-HH control & \checkmark & \checkmark & \checkmark & \checkmark \\
Municipality linear trends & $\times$ & \checkmark & $\times$ & \checkmark \\
\bottomrule
\end{tabular*}\par}
\vspace{4pt}
{\footnotesize
\textit{Notes}: Each column represents a separate regression. The dependent variable is the test score measured in standard deviations. Peabody is the Peabody Test Score and CBCL is the Child Behavior Check List measure, oriented so that higher means fewer problems. Columns 1 and 2 reproduce Columns 3 and 4 of Table 3 in Gillmore (2026): children exposed between in utero and age four, with scores measured seven years later, at ages 7--11; our replication of his Column 3 on the public ELPI files gives $-$0.160 for Peabody and $-$0.103 for CBCL2. Columns 3 and 4 estimate the same two specifications on the 2024 wave: the same exposed children measured fourteen years later, at ages 15--18, against the comparison group conceived after the earthquake and measured in 2017; pooling the two waves adds survey-wave fixed effects, as in Gillmore's short-term table. Standard errors clustered at the municipal level appear in parentheses---number of clusters: 145 in Columns 1 and 2 and 116 in Columns 3 and 4. The estimates use sampling weights. Child/Mother/HH characteristics include gender, age at the moment of the test (in months), birth order, mothers' age and schooling, father present, number of household members, prior mental health condition for mother, father and relative. * p $<$ 0.1, ** p $<$ 0.05, *** p $<$ 0.01.
\par}
\end{minipage}
\end{table}
% ---- end paper/tables/t2_largo_plazo.tex
